\section{House Robber II (Houses in a Circle)}
\label{sec:dp-house-robber-2}

\problemstatement{Same as Section~\ref{sec:dp-house-robber}, except the
houses are arranged in a \textbf{circle} -- house $0$ and house $n-1$ are
now also considered adjacent, so robbing both of them together is
forbidden, in addition to the usual no-two-consecutive-houses rule.}

\begin{patternbox}
\emph{``Same problem, but the arrangement is circular instead of
linear''} is a keyword for a \textbf{reduction} technique: rather than
deriving a new recurrence that somehow accounts for a wraparound
adjacency, we instead run the \emph{already-solved} linear House Robber
(Section~\ref{sec:dp-house-robber}) twice, on two carefully chosen
linear sub-ranges, and combine their results.
\end{patternbox}

\subsection*{Understanding the problem carefully}

The only new constraint, compared to the linear version, is that house
$0$ and house $n-1$ cannot \emph{both} be robbed. Equivalently: in any
valid circular solution, \textbf{at least one} of house $0$ or house
$n-1$ is never robbed. This means every valid circular solution is also
a valid \emph{linear} solution either to the sub-array excluding house
$0$ (i.e.\ houses $1$ through $n-1$), or to the sub-array excluding
house $n-1$ (i.e.\ houses $0$ through $n-2$) -- because whichever of the
two ``forbidden together'' houses is excluded from the solution can
simply be excluded from consideration entirely.

\begin{ideabox}[Reduce to Two Linear Sub-problems, Take the Better One]
Run Section~\ref{sec:dp-house-robber}'s linear House Robber on the
sub-array $\textit{nums}[0..n-2]$ (excluding the last house), and again
on $\textit{nums}[1..n-1]$ (excluding the first house). The answer to the
circular version is the \textbf{maximum} of these two results.
\end{ideabox}

\subsection*{Why considering exactly these two sub-ranges is sufficient}

Every valid circular solution falls into (at least) one of two
categories: those that do not rob house $n-1$ (entirely captured by the
linear sub-problem on houses $0$ through $n-2$, since within that range
the usual linear adjacency rule is the \emph{only} remaining constraint),
and those that do not rob house $0$ (captured by the linear sub-problem
on houses $1$ through $n-1$). A solution that robs \emph{neither} house
$0$ nor house $n-1$ is captured by both sub-problems redundantly, which
costs nothing -- we just take whichever of the two computed maximums is
larger, and that larger value is guaranteed to be a valid circular
answer that misses no better possibility, since any valid circular
solution at all must belong to at least one of these two categories.

\subsection*{Approach 1: Recursion (via Two Calls to the Linear Helper)}

\begin{lstlisting}
#include <bits/stdc++.h>
using namespace std;

int robRecursive(int i, vector<int>& nums) {
    if (i == 0) return nums[0];
    if (i < 0) return 0;

    int skip = robRecursive(i - 1, nums);
    int take = robRecursive(i - 2, nums) + nums[i];
    return max(skip, take);
}

int robLinear(vector<int> sub) {
    if (sub.empty()) return 0;
    return robRecursive((int)sub.size() - 1, sub);
}

int robCircular(vector<int>& nums) {
    int n = nums.size();
    if (n == 1) return nums[0];

    vector<int> excludeLast(nums.begin(), nums.end() - 1);
    vector<int> excludeFirst(nums.begin() + 1, nums.end());

    return max(robLinear(excludeLast), robLinear(excludeFirst));
}

int main() {
    vector<int> nums = {2, 3, 2};
    cout << robCircular(nums) << endl; // 3
    return 0;
}
\end{lstlisting}

\begin{complexitybox}[Approach 1 Complexity]
\textbf{Time.} $O(2^n)$ for each of the two linear calls, so $O(2^n)$
overall (the constant factor of $2$ calls does not change the
asymptotic class).
\textbf{Space.} $O(n)$ recursion stack.
\end{complexitybox}

\subsection*{Approaches 2--4: Memoization, Tabulation, and Space Optimization}

Each of these three later stages is obtained by substituting
Section~\ref{sec:dp-house-robber}'s corresponding memoized, tabulated, or
space-optimized linear House Robber implementation in place of
$\textit{robLinear}$ above -- no new derivation is required, since the
reduction itself (call the linear solver twice, take the max) is
independent of which of the four stages that linear solver internally
uses.

\begin{lstlisting}
#include <bits/stdc++.h>
using namespace std;

int robLinearSpaceOptimized(vector<int>& nums) {
    if (nums.empty()) return 0;
    int n = nums.size();
    int prev2 = 0, prev1 = nums[0];

    for (int i = 1; i < n; i++) {
        int skip = prev1;
        int take = prev2 + nums[i];
        int current = max(skip, take);
        prev2 = prev1;
        prev1 = current;
    }
    return prev1;
}

int robCircularSpaceOptimized(vector<int>& nums) {
    int n = nums.size();
    if (n == 1) return nums[0];

    vector<int> excludeLast(nums.begin(), nums.end() - 1);
    vector<int> excludeFirst(nums.begin() + 1, nums.end());

    return max(robLinearSpaceOptimized(excludeLast),
               robLinearSpaceOptimized(excludeFirst));
}

int main() {
    vector<int> nums = {2, 3, 2};
    cout << robCircularSpaceOptimized(nums) << endl; // 3
    return 0;
}
\end{lstlisting}

\subsection*{Dry run}

\dryrun{Take $\textit{nums} = [2,3,2]$.}

\begin{itemize}
  \item Exclude last (sub-array $[2,3]$): linear House Robber gives
        $\max(2,3)=3$.
  \item Exclude first (sub-array $[3,2]$): linear House Robber gives
        $\max(3,2)=3$.
  \item Overall answer: $\max(3,3)=3$.
\end{itemize}

Checking by hand: robbing houses $0$ and $2$ together is forbidden
(circular adjacency), so the best options are house $1$ alone ($3$), or
house $0$ alone ($2$), or house $2$ alone ($2$) -- confirming $3$ is
correct.

\begin{complexitybox}[Final (Space-Optimized) Complexity]
\textbf{Time.} $O(n)$ -- two linear passes, each $O(n)$.
\textbf{Space.} $O(1)$ extra (beyond the two $O(n)$ sliced sub-arrays,
which could themselves be avoided by passing index ranges instead of
copying, reducing auxiliary space to true $O(1)$).
\end{complexitybox}

\begin{takeawaybox}
When a new constraint (here, circularity) only affects the boundary
between two elements, check whether the problem can be split into a
small, fixed number of cases that each remove the offending boundary
entirely -- reducing the new problem to multiple calls of an
already-solved simpler one, rather than re-deriving a new recurrence
from scratch. This reduction technique is often faster to get correct
than attempting to encode the wraparound constraint directly into a
single, more complex recurrence.
\end{takeawaybox}
